% design-notes.tex -- the engineering reasoning of PCB-0018-A, shared by the
% owner's guide and the design doc. Plain LaTeX only (no house-style macros),
% so both documents can \input it. Facts come from requirements.md,
% architecture/, routing/rules-added.md, reports/*-waivers.json and COMPARISON.md.

\subsection*{Gate drive}

The six power FETs (NTMFS5C628NL, 60 V, about 2 mOhm at 10 V gate drive,
5 x 6 mm) are driven by a TI DRV8300D, a three-phase half-bridge gate driver
with its bootstrap diodes built in. The MCU drives it with six separate PWM
lines (6-PWM mode), and dead time is set in hardware by the resistor on its DT
pin (R301, 40.2 k), so a firmware mistake cannot shoot through a leg on its own.

The driver's gate supply, GVDD, does not come from the motor bus. The LM5164
buck makes +5 V from VM, and a small TPS61040 boost lifts that to +12 V for
GVDD. That is the point of the design: the low-side gates see 12 V and the
high-side gates about 11 V (12 V less the bootstrap diode drop and some droop)
at every bus voltage from 10 V up to the TVS clamp. The requirement was at least
8 V of gate drive at VM = 10 V. We looked at the DRV8320H first, which makes its
gate voltage with a charge pump from VM, but its datasheet guarantees only 6.3 V
high-side drive at VM = 10 V, which fails that floor, so we dropped it.

Each high side runs from a 1 uF bootstrap capacitor (C301--C303) charged through
the driver's internal diode whenever its low side is on. A bootstrap supply
cannot hold a high side on forever, so the firmware must cap the modulation at
about 95\% duty and pulse the low sides once before the first high-side turn-on.
Each gate has a 10 ohm series resistor (R411--R416) to slow the edges and damp
ringing, and a 10 k pull-down to its source so the FETs stay off while the MCU is
in reset or the +12 V rail is still coming up. The driver's own undervoltage
lockouts on GVDD and on each bootstrap rail back that up. The DRV8300D has no
drain-source overcurrent monitor, so overcurrent protection is the MCU's job:
the current-sense outputs also feed the STM32G4's comparators, which are wired to
the motor timer's break input.

\subsection*{Current sensing and the Kelvin connections}

Each leg's low-side FET returns to ground through a 3 mOhm, 3 W, 2512 shunt
(R401--R403). A 3 mOhm shunt drops only 30 mV at 10 A, so a few hundred
microvolts picked up from the ground copper would be a visible error. For that
reason the sense lines do not start from the ground pour. Each pair
(ISNS\_x\_P and ISNS\_x\_N) starts at the shunt's own pads through a net tie,
a deliberate short between two nets that KiCad lets us place exactly, so the
amplifier measures the voltage across the resistor body and not across the
copper that carries 10 A next to it. That is four-wire (Kelvin) sensing.

The amplifiers are INA240A1s, gain 20, with REF1 tied to +3.3 V and REF2 to
ground, so the output sits at 1.65 V at zero current and moves 60 mV per amp.
That gives about $\pm$27 A full scale, and the 20 A peak lands at
$\pm$1.2 V around mid-scale. The INA240 was chosen over the cheaper INA181 because
its reference split gives mid-supply with no extra parts and it tolerates a
common-mode range of $-4$ V to 80 V, which matters when the shunt's ground end
bounces during switching. A small RC filter sits at each amplifier's inputs.

One thing here did not come out as planned. The design checks want the two sense
lines of a pair to run side by side, with no more than 3 mm apart. The INA240's
two inputs sit on opposite rows of its package, 5.4 mm apart, and the shunt's
two net ties are 3.46 mm apart, so each pair has roughly 9 mm where the lines
cannot run together. The packages fix that length, not the routing. The six
findings are waived, and the owner approved that waiver. The INA240's own PWM
rejection and the low bandwidth of the sensed signal make this acceptable, in
our judgement, but it is a judgement, not a measurement.

\subsection*{Power loop}

The loop that matters most in a motor driver is the commutation loop: from the
bus capacitors, through the high-side FET, the low-side FET and the shunt, and
back to the capacitors. Every nanohenry in it turns into voltage overshoot on
the switch node at each edge. Each leg therefore has its own three 10 uF / 50 V
1210 capacitors and a 100 nF placed straight across that leg, from the
high-side drain to the shunt's ground end. The board has four layers: In1 is
an unbroken ground plane under the whole power stage, and In2 carries VM under
it, so the capacitors' return path runs directly under the loop.

The outer layers are 2 oz copper. The high-side drain pads have fields of
12--16 extra vias down to the VM plane, and the shunt ground returns have 12, 12
and 16 vias, after the design checks found the first layout dropping 10 A through
only a few vias per leg. Two 470 uF / 63 V electrolytics hold the bus at low
frequency, and an SMCJ28A TVS clamps transients at about 45 V. Reverse-polarity
protection is a FET (Q101) driven by an LM74500 controller, so current can flow
back into the supply during braking instead of pumping up the bus. There is no
fuse on the board: the supply needs an external 20 A slow fuse.

The bulk capacitors are the weak point for continuous current. At 10 A RMS per
phase about 5--6 A of ripple flows in the bus capacitors, most of it in the
electrolytics at 20 kHz. Each 470 uF can is rated around 1--2 A of ripple (a
guess; we could not get its datasheet), so 10 A continuous is more than the two
cans are rated for. Running the PWM nearer 40 kHz moves more ripple into the
ceramics.

\subsection*{Thermal}

The board has no heatsink, so the copper is the heatsink. Our estimate at 10 A
RMS per phase, 28 V and 25 kHz is about 3.4 W for the whole board: 1.8 W in the
six stage FETs, 0.45 W in the shunts, 0.4 W in the reverse-polarity FET and
0.7 W in the regulators. A 4-layer, 2 oz board in still air sheds roughly
10--12 W per square metre per kelvin over both faces. At 78 x 78 mm that is
about 7--8 C/W, so a mean rise of roughly 25 C with hotter spots at the FETs.
That is a hand estimate scaled from TI's TIDA-010251 reference design, which
measured 49 C at the FETs at 10 A RMS on a board about 60 x 60 mm. No thermal
simulation or measurement has been done for this board, and the 10 A RMS rating
stays provisional. Treat it as the target, not a proven rating.

\subsection*{Size}

The board came out at 78 x 78 mm. The reference board looked roughly 60 x 60
mm in the video frame, but that is our guess from one frame. We put the gate
loops and power loops first and let size follow: three legs of power copper, the
reverse-polarity FET, a pair of 470 uF cans, a 100 V buck and a 12 V boost do not
fit in 60 x 60 mm without crowding the gate loops. The extra area also helps the
thermal estimate above.

\subsection*{What the checks waived, and what that means}

ERC, placement, routed DRC, the design checks (verify) and DFM all pass. The
design checks pass with 390 waivers, every one with its position and reason in
\texttt{reports/verify-waivers.json} and the rule changes behind them in
\texttt{routing/rules-added.md}. They split like this:

\begin{center}
{\small
\begin{tabular}{p{0.22\linewidth} r p{0.58\linewidth}}
\hline
Check & Count & What they are \\
\hline
Current & 359 & 282 via-count, 57 track-width and 20 pour-neck findings. Most (244)
are on GND: stitching vias and small taps into the solid In1 plane that carry
only their own part's return. The rest are Kelvin, gate-return, sense-divider
and snubber taps on the phase and shunt nets, which carry signal current,
and plane edges off the power path. \\
Creepage & 24 & 7 inside the gate driver's (U301) courtyard and 5 inside the
reverse-polarity controller's (U101) courtyard, where the package's own pin
pitch sets the spacing (IPC-2221B Table 6-1, A6, component leads, 0.13 mm up to
100 V); 12 between nets of the same phase leg, which sit within about 15 V of each
other. The owner set this scope on 2026-09-28. \\
Current-sense pairs & 6 & Skew and uncoupled length on ISNS\_A/B/C, set by the
INA240 and shunt pad geometry (above). Owner-approved. \\
Power distribution & 1 & The check reads GND as a supply with no decoupling. GND
is the return, so this is an artifact of the check. \\
\hline
\end{tabular}}
\end{center}

Placement has 14 more waivers of its own (\texttt{reports/place-waivers.json}):
12 where the placement check treats track-width rule areas as keepouts, and two
courtyard overlaps that KiCad's own courtyard check passes.

One fix needs saying plainly. 109 short pad exits on GND, VM and the phase and
shunt nets were narrower than the track-width rule for those nets. They were
turned into solid copper zones of exactly the same shape, which the track-width
rule does not measure. The copper did not get wider. Those necks are no better
than before; they are just no longer counted. They are short pad exits, but if
anything on this board runs hot first, look there.

\subsection*{What has not been done}

Nothing has been built, powered or measured, and nothing has been ordered. Every
current, temperature and ripple figure in this document is a calculation or a
guess. The reverse-polarity FET, the bulk capacitors' ripple rating and the
10 A RMS continuous rating are the three things a first build should check.
